% =========================================================================
\chapter{The project in one page}
\label{ch:overview}
% =========================================================================

\section*{The question}
Can a computer warn, one day ahead, that tomorrow will be an \emph{extreme heat day} at a given place in Himachal
Pradesh? And because such days are rare, does it help to \emph{make extra examples} of them with generative models,
guided by extreme value theory?

\section*{The story in six steps}
\begin{enumerate}
  \item \textbf{Data.} Daily weather for 45 years (1981--2025) on a grid of 529 small squares (cells) covering the
        state, from a data set called ERA5-Land, plus the height and slope of each cell from a terrain map called SRTM.
  \item \textbf{Extreme day.} For each cell and season, the hottest 1\% of days in 1981--2010 sets a threshold. A day
        hotter than that threshold is an extreme heat day.
  \item \textbf{Rows.} Each cell on each day becomes one row of 57 numbers describing what is known today. The answer
        (label) is: is \emph{tomorrow} extreme, yes (1) or no (0)?
  \item \textbf{The difficulty.} Only about 1 row in 100 has the answer ``yes''. A model trained on such data sees
        very few extreme days, and the most severe ones are rarer still.
  \item \textbf{The idea.} Use three generative models (TVAE, CTGAN, a diffusion model) to create extra extreme rows,
        and use extreme value theory (EVT) to decide how severe those extra days should be: mostly the rarest,
        hottest kind, but never hotter than physically possible.
  \item \textbf{The test.} Train the same classifier (LightGBM) with and without the extra rows, and score it on years
        it never saw (2016--2025).
\end{enumerate}

\begin{figure}[H]
\centering
\begin{tikzpicture}[node distance=4mm]
\node[bdat,text width=2.5cm] (d) {ERA5-Land weather\\+ SRTM terrain\\529 cells, 45 years};
\node[blue box,text width=2.5cm,right=of d] (l) {Threshold = top 1\%\\label: is tomorrow\\extreme? (1/0)};
\node[blue box,text width=2.5cm,right=of l] (r) {Rows of\\57 predictors\\(one per cell-day)};
\node[orange box,text width=2.5cm,right=of r] (g) {Generators make\\extra extreme rows,\\EVT sets severity};
\node[green box,text width=4.0cm,below=9mm of l] (c) {LightGBM trained on 1981--2010\\with / without extra rows};
\node[green box,text width=4.0cm,right=8mm of c] (t) {Scored on 2016--2025:\\how well are the\\extreme days found?};
\draw[ar] (d) -- (l); \draw[ar] (l) -- (r); \draw[ar] (r) -- (g);
\draw[ar] (g.south) |- ++(0,-0.45) -| (c.north);
\draw[ar] (c) -- (t);
\end{tikzpicture}
\caption{The whole project in one picture.}
\end{figure}

\section*{The answer}
\begin{itemize}
  \item With real data only, the classifier scores \textbf{0.441} on the most severe days (PR-AUC, explained in
        Chapter~\ref{ch:metrics}; 1 is perfect).
  \item Adding extreme days from TVAE or CTGAN raises this to \textbf{0.601 (+36\%)}; from the diffusion model to
        \textbf{0.579 (+31\%)}.
  \item EVT is the reason it works: putting most added days in the rarest part of the tail improved every generator
        (by up to 14\%), while adding days without that focus helped much less.
\end{itemize}

\begin{sayit}
``Extreme heat days are 1\% of the data, so a model rarely sees them. I used generative models to create more of them,
and extreme value theory to make most of the new days the rare, severe kind, within physical limits. On years the
model never saw, this made its warnings for the most dangerous days 31 to 36\% better.''
\end{sayit}
